% Generated by scripts/assemble.py; corrections belong in source/corrections.json.
\examhead{75}{吉林大学 2024 年数学分析试卷}{74}
% p075-b01; source PDF page 75
\par\Needspace{5\baselineskip}
\blockanchor{p075-b01}
\majorquestion{一}求极限.\par
% p075-b02; source PDF page 75
\blockanchor{p075-b02}
\question{1}$\lim_{n\to\infty}\frac{1+\frac{1}{2}+\cdots+\frac{1}{n^{2}-1}+\frac{1}{n^{2}}}{n\ln n}.$\par
% p075-b03; source PDF page 75
\blockanchor{p075-b03}
\question{2}$\lim_{n\to\infty}n^{2}\left(n\sin\frac{1}{n}-1\right).$\par
% p075-b04; source PDF page 75
\blockanchor{p075-b04}
\question{3}$\lim_{x\to 0}\frac{\sqrt[3]{1 + \tan^2x} - \sqrt{\cos x}}{x\sin x}.$\par
% p075-b05; source PDF page 75
\blockanchor{p075-b05}
\question{4}$\lim_{n\to \infty}n\int_0^1 x^n\cos (1 + x^2)\mathrm{d}x.$\par
% p075-b06; source PDF page 75
\blockanchor{p075-b06}
\question{5}\displayformula{\lim _ {n \rightarrow \infty} \frac {1}{n} \left(\frac {1}{\sqrt {n ^ {2} + 1}} + \frac {1}{\sqrt {n ^ {2} + 2}} + \dots + \frac {1}{\sqrt {2 n ^ {2} - 1}} + \frac {1}{\sqrt {2 n ^ {2}}}\right)}\par
% p075-b07; source PDF page 75
\par\Needspace{6\baselineskip}
\blockanchor{p075-b07}
\question{6}设 $u = u(x, y), v = v(x, y)$ 是由方程组\par
\nopagebreak[4]
% p075-b08; source PDF page 75
\blockanchor{p075-b08}
\displayformula{\left\{ \begin{array}{l} x u - y v = 0; \\ y u + x v = 2. \end{array} \right.}
% p075-b09; source PDF page 75
\blockanchor{p075-b09}
在点 $(x_0, y_0, u_0, v_0) = (1, 1, 1, 1)$ 的某个邻域内定义的隐函数，求 $\frac{\partial u}{\partial x}(1, 1), \frac{\partial u}{\partial y}(1, 1), \frac{\partial v}{\partial x}(1, 1), \frac{\partial v}{\partial y}(1, 1)$.\par
% p075-b10; source PDF page 75
\blockanchor{p075-b10}
\question{7}求曲面 $2z = x^{2} + y^{2}$ 与 $z = \sqrt{x^2 + y^2}$ 围成区域的体积.\par
% p075-b11; source PDF page 75
\blockanchor{p075-b11}
\question{8}求第二型曲线积分 $\int_{L} \frac{x \, \mathrm{d}y - y \, \mathrm{d}x}{2x^{2} + y^{2}}$, 其中 $L$ 是圆周 $x^{2} + y^{2} = 1$, 方向为逆时针.\par
% p075-b12; source PDF page 75
\par\Needspace{5\baselineskip}
\blockanchor{p075-b12}
\majorquestion{二}求 $f(x, y) = 2x^{2} + 4xy - 2y^{2}$ 在闭区域 $D = \{(x, y) \mid x^{2} + y^{2} \leq 5\}$ 上的最大值和最小值.\par
% p075-b13; source PDF page 75
\par\Needspace{5\baselineskip}
\blockanchor{p075-b13}
\majorquestion{三}求级数 $\sum_{n=1}^{\infty} \frac{(-1)^{n-1}}{n2^n}$ 的和.\par
% p075-b14; source PDF page 75
\par\Needspace{5\baselineskip}
\blockanchor{p075-b14}
\majorquestion{四}设  $a_{n}=\left(1-\frac{2\ln(\ln n)}{n}\right)^{n}, n=1,2,\cdots$ ，判断级数  $\sum_{n=1}^{\infty}a_{n}$  的敛散性.\par
% p075-b15; source PDF page 75
\par\Needspace{5\baselineskip}
\blockanchor{p075-b15}
\majorquestion{五}证明: $\lim_{h \to 0^{+}} \int_{0}^{1} \frac{h}{h^{2} + x^{2}} \ln (x^{2} + 2) \, \mathrm{d}x = \frac{\pi}{2} \ln 2.$\par
% p075-b16; source PDF page 75
\par\Needspace{5\baselineskip}
\par\Needspace{6\baselineskip}
\blockanchor{p075-b16}
\majorquestion{六}设函数 $f(x)$ 在 $[0,2]$ 上二阶导数连续，且 $f(0) = f(2) = 0$，证明：\par
\nopagebreak[4]
% p075-b17; source PDF page 75
\blockanchor{p075-b17}
\displayformula{\int_ {0} ^ {2} f (x) \mathrm{d} x \leq \frac {2}{3} \max _ {x \in [ 0, 2 ]} \{| f ^ {\prime \prime} (x) | \}.}
% p075-b18; source PDF page 75
\par\Needspace{5\baselineskip}
\blockanchor{p075-b18}
\majorquestion{七}设 $F$ 为定义在 $\mathbb{R}^n$ 上的函数.\par
% p075-b19; source PDF page 75
\blockanchor{p075-b19}
(1) 函数 $F$ 不一致连续的充要必要条件是: 存在点列 $\{\pmb{x}^{(n)}\}$ 和 $\{\pmb{y}^{(n)}\} \subset \mathbb{R}^n$ 和常数 $a$, 其中有 $\lim_{n \to \infty} \| \pmb{x}^{(n)} - \pmb{y}^{(n)} \| = 0$, 但对 $n = 1, 2, \dots$, 都有 $|F(\pmb{x}^{(n)}) - F(\pmb{y}^{(n)}) \geq a$.\par
% p075-b20; source PDF page 75
\blockanchor{p075-b20}
(2) 若  $F(x_{1}, x_{2}, \cdots, x_{n}) = f(x_{1}^{2} + x_{2}^{2} + \cdots + x_{n}^{2})$ ，f 在  $[0, +\infty)$  上连续可导，同时  $\lim_{t \to +\infty} f'(t) = b \neq 0$ ，证明：F 在  $R^{n}$  上不一致连续.\par
